
\DocumentMetadata{tagging=on,
	tagging-setup={math/setup=mathml-SE},
	pdfstandard=ua-2,
	lang=en-US
}
\documentclass{article}

%Math Libraries
\usepackage{multirow, longtable, amsmath, amssymb, amsthm}
\usepackage[mathscr]{euscript}

%Formatting
\usepackage[margin = 1 in]{geometry}
\usepackage{indentfirst}
\usepackage{graphicx}
\usepackage{fancyhdr}
\usepackage{tabularx}
\usepackage{cite}

%Diagram Packages
\usepackage{tikz}
\usetikzlibrary{automata,positioning,arrows}

%Standard Commands
\newcommand{\mm}{\mathscr}
\newcommand{\B}{{\mathcal{B}}}
\newcommand{\A}{{\mathcal{A}}}
\newcommand{\N}{{\mathbb{N}}}
\newcommand{\R}{{\mathbb{R}}}
\newcommand{\C}{{\mathbb{C}}}
\newcommand{\Q}{{\mathbb{Q}}}
\newcommand{\Z}{{\mathbb{Z}}}


\newcommand{\abs}[1]{{|{#1}|}}
\newcommand{ \norm }[1]{{ || {#1}||}}
\newcommand{\cl}[1]{{\overline{#1}}}
\newcommand{\pd}{\partial}
\newcommand{\ip}[2]{ \langle {#1},{#2}\rangle }
\newcommand{\ls}{\lesssim}
\newcommand{\gs}{\gtrsim}
\newcommand{\supp}{\text{supp}}

%Result Environment
\newcounter{result}[section]
\newenvironment{res}[1][]{\refstepcounter{result}\par\medskip
	\noindent \textbf{[\thesection.\theresult] #1} \rmfamily}{\medskip}

\newenvironment{defn}[1][]{\refstepcounter{result}\par\medskip \noindent \textbf{Definition [\thesection.\theresult] #1} \rmfamily}{\medskip}

%Proof Environment
\newenvironment{pf}{%
	\par%
	\medskip
	\leftskip=2em\rightskip=2em%
	\noindent\ignorespaces \textit{Proof:} \newline}{%
	\,\, \newline \textit{Q.E.D.}\par\medskip}

\title{Math 126 Summer 2026 HW 3}
\author{Lec 001, Ebbighausen}


\begin{document}
	\pagestyle{fancy}
	\fancyfoot{} 
	\fancyfoot[L]{Math 126}
	
	\maketitle
	
	\textbf{Problem 1:} (Borthwick 7.6) On HW 3, we investigated a quantity
	\begin{align*}
		\eta(t) = \int_{\Omega} u(t,x)^{2} dx
	\end{align*}
	for $u$ a smooth solution to the heat equation on $(0,\infty) \times \Omega$ with bounded $C^{1}$ domain $\Omega$. Suppose $u$ is such a solution and $u|_{t=T} = 0$, $u|_{\pd \Omega} = 0$. 
	
	Part A.) Use Cauchy-Schwarz to deduce that 
	\begin{align*}
		\eta'(t)^{2} \leq 4 \eta(t) \int_{\Omega} \lvert \frac{\pd u}{\pd t} \rvert^{2} dx
	\end{align*}
	
	Part B.) Show 
	\begin{align*}
		\eta''(t) = 4 \int_{\Omega} \lvert \frac{\pd u}{\pd t} \rvert^{2} dx
	\end{align*}
	
	so that from A.), $\eta'(t)^{2} \leq \eta(t) \eta''(t)$. 
	
	
	Part C.) Suppose $\eta(0) > 0$ and $\eta$ is defined in some neighborhood of $0$. Using part B.), show 
	\begin{align*}
		(\log(\eta(t)) '' \geq 0
	\end{align*}
	
	This implies that $\log(\eta(t))$ is \textit{convex} and, in particular, bounded below by its tangent lines such that 
	\begin{align*}
		\log(\eta(t)) \geq \log(\eta(0)) + \frac{\eta'(0)}{\eta(0)}t
	\end{align*}
	so that $\eta(t) \geq \eta(0) e^{-ct}$. Thus, if $\eta(0)>0$, $\eta$ is strictly positive for all time. 
	
	Part D.) Conclude from Part C.) that if $\eta(T) = 0$, $\eta(t) = 0$ for all $t$ and $u=0$ (it may help to reference HW 3). 
	
	\textbf{Problem 2:} In the separation of variables section, we used the Legendre polynomials to help construct solutions. These are an orthonormal basis for $L^{2}$ via polynomials, as opposed to exponentials. 
	
	Part A.) Use Gram-Schmidt on $\{1,x,x^{2}\}$ with the $L^{2}((-1,1))$-inner product to obtain the first 3 Legendre polynomials. Verify that these are normalized (norm 1) versions of the first three solutions of 
	\begin{align*}
		P_{k}(x) = \frac{1}{2^{k}k!} \frac{d^{k}}{dx^{k}}(x^{2}-1)^{k}
	\end{align*}
	
	
	Part B.) These polynomials are solutions of the ODE
	
	\begin{align*}
		\frac{d}{dx} \left[ (x^{2}-1) \frac{d}{dx} f \right] -k(k+1)f = 0
	\end{align*}
	
	or eigenvectors $Lf = k(k+1)f$ for $L = \frac{d}{dx} \left[ (x^{2}-1) \frac{d}{dx} \right]$. For $u,v \in C^{2}([-1,1])$, check that 
	\begin{align*}
		\langle u,Lv \rangle = \langle Lu,v\rangle   
	\end{align*}
	so that the $P_{k}$ with distinct $k$ are orthogonal. 
	
	
	Part C.) Let $f(x) \in L^{2}((-1,1))$ be a function so 
	\begin{align*}
		\int_{-1}^{1} f(x)x^{l}dx = 0
	\end{align*}
	for $l=0,1,2,...$. Show that 
	\begin{align*}
		\int_{-1}^{1} f(x)S_m(x)dx = 0
	\end{align*}
	
	for $S_{m}(x) = \sum_{l=0}^{m} \frac{(-ikx)^{l}}{l!}$. Since $\lim_{m \to \infty} S_{m}(x) = e^{-ikx}$, show $f$ must be $0$. 
	
	Part D.) Use Part C.) to show the Legendre polynomial $P_{k}$ are a basis for $L^{2}$. 
	
	

	
\end{document}
